
\documentclass{article}
\usepackage{colortbl}
\usepackage{makecell}
\usepackage{multirow}
\usepackage{supertabular}

\begin{document}

\newcounter{utterance}

\centering \large Interaction Transcript for game `ifeval', experiment `ifeval', episode 339 with selfrefine{-}Qwen3.5{-}9B.
\vspace{24pt}

{ \footnotesize  \setcounter{utterance}{1}
\setlength{\tabcolsep}{0pt}
\begin{supertabular}{c@{$\;$}|p{.15\linewidth}@{}p{.15\linewidth}p{.15\linewidth}p{.15\linewidth}p{.15\linewidth}p{.15\linewidth}}
    \# & \multicolumn{2}{c}{Player} && \multicolumn{2}{c}{Game Master} \\
    \hline

    \theutterance \stepcounter{utterance}  
    & & & \multicolumn{4}{p{0.6\linewidth}}{
        \cellcolor[rgb]{0.9,0.9,0.9}{
            \makecell[{{p{\linewidth}}}]{
                \texttt{\tiny{[P1$\langle$GM]}}
                \texttt{What is inside Shinto shrines? Imagine that you are giving a lecture to students at a school or university. Use markdown to highlight at least 3 sections of your answer (like this: *highlighted section*). Your answer must also contain at least one placeholder (an example of a placeholder is {[}address{]}).} \\
            }
        }
    }
     \\ \\

    \theutterance \stepcounter{utterance}  
    & \multicolumn{4}{p{0.6\linewidth}}{
        \cellcolor[rgb]{0.561,0.737,0.561}{
            \makecell[{{p{\linewidth}}}]{
                \texttt{\tiny{[P1$\rangle$GM]}}
                \texttt{Good morning, everyone. Today, we are going to explore the spiritual heart of Japan by looking inside a Shinto shrine. Unlike many other religious traditions where the focus is on a statue of a deity or a written scripture, the interior of a Shinto shrine is designed to facilitate a direct connection between the human world and the *kami*—the spirits or divine forces that inhabit nature, ancestors, and phenomena.} \\
\\ 
\texttt{*The most sacred space you will find inside any shrine is the Haiden (Hall of Prayer) and the Honden (Main Hall).* While the Haiden is where worshippers gather to perform rituals and offer prayers, it is actually the Honden that holds the true sanctity. However, you will rarely see an idol or a statue inside the Honden. Instead, the *kami* is believed to reside in the very structure of the building itself, often symbolized by a sacred mirror, sword, or jewel known as the *Sanshu no Kami*. In many cases, the *kami* is simply present in the air, the wood, and the space, making the building a living vessel for the divine rather than a container for an object.} \\
\\ 
\texttt{*Another essential element you will encounter is the Torii gate and the approach to the shrine.* Before even entering the main structures, visitors pass through a Torii gate, which marks the transition from the mundane world to the sacred space. Inside, you will often see a temizuya (water purification fountain) where visitors cleanse their hands and mouth before approaching the altar. This act symbolizes the purification of the body and spirit, ensuring that no impurities disturb the sacred atmosphere within. The layout is designed to guide you through a journey of increasing reverence, culminating in the silent, awe{-}inspiring presence of the *kami*.} \\
\\ 
\texttt{To understand the specific location and history of the shrine we are discussing today, please refer to the detailed map located at {[}address{]}, which outlines the architectural layout and the specific *kami* enshrined there. By observing these elements, we can see that Shinto shrines are not about worshipping images, but about honoring the sanctity of the space and the invisible forces that sustain our world.} \\
\\ 
\texttt{Thank you for your attention, and I look forward to your questions.} \\
            }
        }
    }
    & & \\ \\

    \theutterance \stepcounter{utterance}  
    & & & \multicolumn{2}{p{0.3\linewidth}}{
        \cellcolor[rgb]{0.9,0.9,0.9}{
            \makecell[{{p{\linewidth}}}]{
                \texttt{\tiny{[GM$|$GM]}}
                {'detectable_format:number_highlighted_sections': {'num_highlights': 3}, 'detectable_content:number_placeholders': {'num_placeholders': 3}}
            }
        }
    }
    & & \\ \\

    \theutterance \stepcounter{utterance}  
    & & & \multicolumn{2}{p{0.3\linewidth}}{
        \cellcolor[rgb]{0.9,0.9,0.9}{
            \makecell[{{p{\linewidth}}}]{
                \texttt{\tiny{[GM$|$GM]}}
                \texttt{game\_result = LOSE} \\
            }
        }
    }
    & & \\ \\

\end{supertabular}
}

\end{document}
